% Section 11.1, from lectures/L11/appendix/a_frameworks.md.

\section{Kernel Frameworks}
\label{c11:app:a}

Why your driver probably should not be a character device, what a subsystem gives you in exchange
for implementing its interface, and how to tell which one a device belongs to.

The argument in one sentence: \textbf{a character device is an interface you invent and then
maintain forever; a framework is an interface that already exists and that other people's tools
already speak.}

\subsection{What you have, and what you do not}
\label{c11:app:a1}

By Chapter~\ref{c10:ch} the driver works. It probes off the device tree and services its
interrupt, and in Chapter~\ref{c9:ch} it presented \file{/dev/qa_wait} with a \code{read()} you
wrote. Nothing is broken.

Now list what it does not have:
\begin{itemize}
  \item No way for a program that does not know about it to find it.
  \item No documented interface. Yours exists in your head and in the source.
  \item No existing tools. Every user writes their own reader.
  \item No power management. Suspend and resume do nothing.
  \item No way to express what the device \emph{is}, only what it does.
\end{itemize}
Each of those is something a subsystem provides, and provides identically for every driver in it.
That is the trade: \textbf{you implement an interface someone else designed, and you stop having to
design, document, version and support one of your own.}

\lecturefigure{lectures/L11/appendix/images/framework_stack.png}
  {Two columns for the same device. A character device of your own needs open and release, read and
   write, an ioctl ABI of your own, a /dev node, a class and a device, and still nothing can find
   it. An IIO device needs a channel specification and read\_raw, and the framework supplies the
   sysfs path, a documented ABI, existing tools and buffered modes.}
  {fig:framework_stack}

\subsection{The cost of the shortcut}
\label{c11:app:a2}

Adding an \code{ioctl} is always the path of least resistance. It is worth knowing what it costs
before taking it.

\textbf{It is an ABI the moment a second program uses it.} The numbers and the structure layout are
frozen; you may add but never change or remove. There is no deprecation period for something a
customer's binary calls.

\textbf{It is a portability trap.} A 32-bit process on a 64-bit kernel passes a differently laid-out
structure, and handling that needs \code{compat_ioctl} and a structure with no implicit padding and
no \code{long}. Getting this wrong is a standing source of kernel bugs, and it is invisible until
somebody runs a 32-bit binary.

\textbf{It is undiscoverable.} \code{ls /sys/bus/iio/devices} finds every IIO device on a machine.
Nothing finds every device that implements your ioctl.

\textbf{And there is usually already a subsystem for it.} The kernel has spent thirty years
accumulating interfaces for classes of hardware, and ``a thing that produces samples'' or ``a thing
with controllable pins'' are both very well-trodden.

\subsection{The tour}
\label{c11:app:a3}

Ten subsystems worth recognising, with the question each answers.

\begin{booktable}{@{}l>{\hsize=0.92\hsize}L>{\hsize=1.08\hsize}L@{}}
  \thead{Subsystem} & \thead{For} & \thead{Userspace sees} \\
  \midrule
  \textbf{gpio} & Individually controllable pins & \code{/dev/gpiochipN}, libgpiod \\
  \textbf{pinctrl} & What a pin is multiplexed \emph{as} & Nothing directly; other drivers \\
  \textbf{clk} & Clock sources and gates & Nothing directly; \code{/sys/kernel/debug/clk} \\
  \textbf{regulator} & Power supplies & Read-only \code{/sys/class/regulator} \\
  \textbf{i2c}, \textbf{spi} & Devices on those buses & \code{/dev/i2c-N}, \code{/dev/spidevN.M} \\
  \textbf{regmap} & Register access over any of those & Nothing; it is a driver-side helper \\
  \textbf{iio} & Anything that samples or converts & \code{/sys/bus/iio/devices/}, \code{/dev/iio:deviceN} \\
  \textbf{input} & Keys, buttons, touch, motion & \code{/dev/input/eventN} \\
  \textbf{mtd} & Raw flash & \code{/dev/mtdN}, \code{/dev/mtdblockN} \\
  \textbf{watchdog} & A timer that reboots the machine & \code{/dev/watchdog} \\
\end{booktable}

\noindent Four of them are worth more than a row.

\textbf{pinctrl} is the one people are surprised to need. Before a pin can be a GPIO it has to be
\emph{multiplexed} as one rather than as a UART line or an I2C signal, and that is a different piece
of hardware in the SoC. Your driver should not be doing it: the pin configuration belongs in the
device tree, and pinctrl applies it before your \code{probe} runs.

\textbf{clk and regulator} exist because the answer to ``turn the device on'' is shared and
refcounted. Two drivers on one clock is the normal case, and a driver that writes the gate register
directly works until the second one disagrees with it. \code{devm_clk_get} and
\code{clk_prepare_enable} cost two lines and are correct with any number of consumers.

\textbf{regmap} is a driver-side helper rather than a userspace interface. Most drivers for an I2C
or SPI chip contain the same fifty lines of read-modify-write, byte-order and caching; regmap
replaces them with a \emph{description}, which registers exist, which are volatile, which are
cacheable, and handles the rest. It is worth reaching for the moment a driver has more than a
handful of registers.

\textbf{iio} is where anything that measures belongs. Not just ADCs: accelerometers, light sensors,
pressure, temperature, DACs, anything with a value and a unit. It carries scale, offset and units so
that userspace can convert raw counts into volts without knowing anything about the chip.

\textbf{Two of the ten you will drive yourself and eight you will only have read about}, which is
worth knowing before the table starts to feel like competence: \S\ref{c11:app:a4} and
\S\ref{c11:app:a5} register \code{qa-dev} with iio and gpio and the lab checks both, but
\code{qa-dev} sits on no bus and has no clock or regulator to ask for, so nothing here can teach you
what regmap feels like against a real I2C part.

\subsection{Registering with IIO}
\label{c11:app:a4}

The whole of it, for a device with one channel:

\begin{ccode}
indio_dev = devm_iio_device_alloc(&pdev->dev, sizeof(*priv));
priv      = iio_priv(indio_dev);

indio_dev->name          = "qa_dev";
indio_dev->info          = &qa_iio_info;      /* .read_raw = qa_read_raw */
indio_dev->modes         = INDIO_DIRECT_MODE;
indio_dev->channels      = qa_channels;
indio_dev->num_channels  = ARRAY_SIZE(qa_channels);

ret = devm_iio_device_register(&pdev->dev, indio_dev);
\end{ccode}

\noindent \code{devm_iio_device_alloc} allocates the IIO device \emph{and} your private structure in
one object; \code{iio_priv()} gets yours back out. That is the framework's idiom for per-device
state and it replaces the \code{devm_kzalloc} from Chapter~\ref{c10:ch}.

The channel says what the value \emph{is}:

\begin{ccode}
static const struct iio_chan_spec qa_channels[] = {
        {
                .type                = IIO_VOLTAGE,
                .indexed             = 1,
                .channel             = 0,
                .info_mask_separate  = BIT(IIO_CHAN_INFO_RAW),
        },
};
\end{ccode}

\noindent and \code{.info_mask_separate} is the list of attributes to create.
\code{IIO_CHAN_INFO_RAW} makes \code{in_voltage0_raw}; adding \code{IIO_CHAN_INFO_SCALE} would make
\code{in_voltage0_scale}, and userspace multiplies the two to get millivolts, the unit the IIO ABI
documents for a voltage.

\code{read_raw} is the one callback:

\begin{ccode}
static int qa_read_raw(struct iio_dev *indio_dev, struct iio_chan_spec const *chan,
                       int *val, int *val2, long mask)
{
        struct qa_priv *priv = iio_priv(indio_dev);

        switch (mask) {
        case IIO_CHAN_INFO_RAW:
                *val = atomic_read(&priv->last_sample);
                return IIO_VAL_INT;
        default:
                return -EINVAL;
        }
}
\end{ccode}

\noindent \textbf{What that buys}, on the target, immediately:

\begin{console}
# ls /sys/bus/iio/devices/
iio:device0
# cat /sys/bus/iio/devices/iio:device0/name
qa_dev
# cat /sys/bus/iio/devices/iio:device0/in_voltage0_raw
195
\end{console}

\noindent No \file{/dev} node was written, no \code{read()} implemented, no \code{ioctl} invented.
And because the device is in counter mode, reading the attribute twice measures the sample rate:

\begin{console}
195 ... one second ... 1076
\end{console}

\noindent 881 samples in a second, which is Chapter~\ref{c8:ch}'s interrupt rate arriving through an
interface that knows nothing about interrupts.

\subsection{Registering a gpiochip}
\label{c11:app:a5}

Five callbacks and a structure:

\begin{ccode}
priv->gc.label           = dev_name(&pdev->dev);
priv->gc.parent          = &pdev->dev;
priv->gc.owner           = THIS_MODULE;
priv->gc.base            = -1;              /* let the core allocate the numbers */
priv->gc.ngpio           = 8;
priv->gc.get             = qa_gpio_get;
priv->gc.set             = qa_gpio_set;
priv->gc.direction_input = qa_gpio_direction_input;
priv->gc.direction_output = qa_gpio_direction_output;
priv->gc.get_direction   = qa_gpio_get_direction;

ret = devm_gpiochip_add_data(&pdev->dev, &priv->gc, priv);
\end{ccode}

\noindent \code{.base = -1} matters: a hardcoded global GPIO number is a legacy practice. This
kernel still accepts one, with a warning that static allocation is deprecated, but two chips whose
hardcoded ranges overlap cannot both register. \code{gpiochip_get_data(gc)} retrieves the pointer
you passed.

Note that \code{.set} returns \code{void} in this kernel. It gained an \code{int} return later, so a
driver copied from newer material will not compile here, and one copied from older material will not
compile there. \S\ref{c3:app:a2}'s unstable internal API, again.

What appears:

\begin{console}
# ls -l /sys/bus/gpio/devices/gpiochip1/of_node
of_node -> .../firmware/devicetree/base/platform-bus@c000000/qa-dev@0
# ls -l /dev/gpiochip1
crw------- 1 root root 254, 1 /dev/gpiochip1
\end{console}

\noindent The chip is a child of your platform device, linked back to the device tree node it came
from, with a character device the framework created and will remove when the driver unbinds.

\textbf{There is no \file{/sys/class/gpio}.} That interface is deprecated and this course's kernel
does not enable it; the character device replaced it, and a driver written against the sysfs
interface is a driver written against something on its way out. \file{kernel/qa.config} says so in a
comment.

\subsection{What userspace can now do}
\label{c11:app:a6}

The tester shipped with the lab contains \textbf{nothing specific to this driver}. It finds the chip
by the label the driver gave it and then speaks the standard ABI:

\begin{ccode}
ioctl(fd, GPIO_GET_CHIPINFO_IOCTL, &info);        /* what is this chip */
ioctl(chip, GPIO_V2_GET_LINE_IOCTL, &req);        /* give me line 0 as an output */
ioctl(line, GPIO_V2_LINE_SET_VALUES_IOCTL, &v);   /* drive it */
\end{ccode}

\begin{console}
found /dev/gpiochip1: label "c000000.qa-dev", 8 lines
  ok  line 0 can be driven high
  ok  line 4 reads high, through the loopback
\end{console}

\noindent That is the argument, demonstrated: \textbf{a program written before your driver existed
can drive your hardware}, because both sides implement an interface neither of them owns.
\code{libgpiod} and its \code{gpioinfo}, \code{gpioget} and \code{gpioset} tools do exactly this and
would work here unchanged.

\subsection{Choosing a subsystem}
\label{c11:app:a7}

The question is \textbf{what the device is}, not what it does.

\begin{booktable}{@{}>{\hsize=1.13\hsize}L>{\hsize=0.87\hsize}L@{}}
  \thead{The device is} & \thead{Subsystem} \\
  \midrule
  Something that measures or converts a value & iio \\
  Something with individually controlled pins & gpio \\
  A source of human input & input \\
  Raw flash & mtd \\
  A timer that resets the board & watchdog \\
  A clock source & clk \\
  A power supply & regulator \\
  On an I2C or SPI bus & i2c or spi, plus one of the above \\
\end{booktable}

\noindent A device is often several: a sensor on I2C that reports temperature is an \textbf{i2c}
client and an \textbf{iio} device, and those are not competing answers.

\textbf{When nothing fits}, a character device is a legitimate answer and not a failure. Write it
with \code{misc_register} rather than the long way, document the interface, use
\code{_IOR}/\code{_IOW} properly, and be honest that you now own an ABI. What is not legitimate is
reaching for it because implementing an interface looked like more work than inventing one.

\subsection{What a framework does not do}
\label{c11:app:a8}

Worth stating, because the argument above is one-sided by design.

\textbf{It does not make your driver correct.} Everything from Chapters~\ref{c6:ch} to~\ref{c9:ch}
still applies: the locking, the interrupt acknowledgement, the sleeping rules. A framework changes
what your driver presents, not what it has to get right underneath.

\textbf{It constrains you to its model.} If your device genuinely does not fit IIO's idea of a
channel, forcing it in produces a driver that is harder to read than a character device would have
been. The model is a good fit surprisingly often and is not a universal one.

\textbf{It can be slower, and here is how much.} Going through a subsystem's generic path costs
something against a purpose-built \code{read()}. Measured on this target, reading
\code{in_voltage0_raw} repeatedly:

\begin{narrowtable}{@{}lll@{}}
  \thead{From} & \thead{Per read} & \thead{Rate} \\
  \midrule
  A C program & 571 us & 1,748 reads/s \\
  A shell loop of \code{cat} & 29,700 us & 34 reads/s \\
\end{narrowtable}

\noindent The shell figure is fifty times worse because each \code{cat} is a fork and an exec, and
it is worth knowing before you write a monitoring script: \textbf{a shell reading sysfs cannot keep
up with a device sampling at more than about thirty hertz.}

Even the C figure is only twice this device's 881 samples a second, so reading the attribute in a
loop is already marginal at one kilohertz and hopeless above it. That is not a criticism of sysfs;
it is what sysfs is for, which is configuration and occasional inspection. \textbf{IIO's buffered
and triggered modes exist for exactly this}, delivering batches through \file{/dev/iio:deviceN}
instead, and a driver that needs them says so by allocating a buffer rather than by abandoning the
framework.

\textbf{And it does not remove the ABI problem, it moves it to someone competent.} The GPIO
character device is on its second version, \code{GPIO_V2_*}, because the first got some things
wrong. The difference is that the migration was designed, documented and supported by the subsystem
maintainers rather than by you.

\lecturefigure[0.68]{lectures/L11/appendix/images/sysfs_cost.png}
  {A logarithmic bar chart of reads per second: 1,748 from a C program at 571 microseconds each,
   against 34 from a shell loop of cat at 29,700 microseconds each, with a dashed line marking the
   881 samples a second the device produces.}
  {fig:sysfs_cost}
